% AgentRoute — EMNLP 2026 Industry Track submission
%
% Double-blind review: no author names, no affiliations, no self-identifying URLs.
% Page limit: 6 pages of main body (references, Limitations, Ethical Considerations,
% appendices, and acknowledgments do not count).
%
% This file uses the standard ACL 2026 style files (acl.sty + acl_natbib.bst).
% Download from https://github.com/acl-org/acl-style-files and place
% acl.sty + acl_natbib.bst alongside this file before compiling.

\documentclass[11pt]{article}

% Use the [review] option for the anonymous submission version.
% For the camera-ready version, replace with \usepackage{acl}.
\usepackage[review]{acl}

\usepackage{times}
\usepackage{latexsym}
\usepackage[T1]{fontenc}
\usepackage[utf8]{inputenc}
\usepackage{microtype}
\usepackage{graphicx}
\usepackage{booktabs}
\usepackage{amsmath}
\usepackage{amssymb}
\usepackage{url}

% Avoid duplicate citation/hyperref warnings; natbib is loaded by acl.sty.

\title{AgentRoute: A Production Communication Layer for \\
       Multi-Agent LLM Systems with Zero Routing Overhead}

\author{Anonymous Submission}

\begin{document}
\maketitle

\begin{abstract}
Large language model (LLM)-based multi-agent systems are increasingly deployed
in production, but their communication overhead grows with both the number of
registered agents and the complexity of the routing policy. Recent routing
methods (MasRouter, RopMura, OI-MAS, RCR-Router, EvoMAS, DyTopo) reduce
message-payload tokens between agents, yet the routing decision itself
typically requires multiple LLM calls, a trained controller, or both---a fixed
cost that becomes the binding constraint at production scale. We present
AgentRoute, a deterministic communication layer that combines a pattern-based
domain classifier, an LRU result cache, and an optional small-model fallback
for low-confidence queries. AgentRoute issues exactly one LLM call to the
chosen specialist per resolved query, uses no learned components, and has
zero training cost. We compare AgentRoute to five 2025--2026 baselines on
MedQA-USMLE and GPQA Diamond using a shared Qwen 2.5--3B backbone, study
scalability up to 200 agents, and report a production case study against the
Claude API. Our results characterise the regime where the routing layer's
fixed cost dominates and a deterministic router is the appropriate design
choice.
\end{abstract}

\section{Introduction}
\label{sec:intro}

Multi-agent LLM systems are now deployed in production across coding
assistants, customer support, agentic workflow automation, and biomedical
question answering. These systems decompose tasks across role-typed agents
that exchange natural-language messages \citep{wu2023autogen,hong2024metagpt,
li2023camel,qian2024chatdev,chen2024agentverse}, and as the number of
registered agents grows, the communication layer between them becomes a
non-trivial fraction of the total system cost. An empirical study of seven
multi-agent frameworks attributes a substantial fraction of system failures to
\emph{inter-agent communication} faults rather than to any individual agent's
reasoning errors~\citep{cemri2025why}, motivating recent work on more
selective routing.

The recent routing literature has converged on increasingly expressive
controllers. MasRouter trains a cascaded controller that jointly chooses
collaboration mode, role assignment, and LLM
backbone~\citep{yue2025masrouter}. RopMura introduces a four-submodule planner
(question splitter, question selector, judger, defender) with
cluster-embedding routing across knowledge agents~\citep{wang2025ropmura}.
OI-MAS uses state-dependent role and model routing with confidence-aware
early termination~\citep{wang2026oimas}. RCR-Router introduces role-aware
context routing under per-role token budgets and an importance
scorer~\citep{liu2025rcrrouter}. EvoMAS maintains and mutates a pool of
multi-agent system configurations~\citep{evomas2026}, and DyTopo rebuilds the
communication topology each round through semantic key/query
matching~\citep{lu2026dytopo}. These methods reduce the \emph{message payload}
between agents, often by 30--70\%, on a range of multi-hop reasoning
benchmarks.

\paragraph{The production gap.}
In production deployments operating at scale, however, what matters is not
only the message payload but the cost of the routing layer itself. A router
that calls an LLM two to three times per query, or that requires
training on per-task data, has a fixed cost that grows with traffic
independent of how much it shrinks the message payload. For the production
operator, three properties of the routing layer dominate: predictable
sub-call latency, bounded per-query token budget, and an auditable decision
trace. None of these are first-class concerns in the recent routing
literature, which is largely benchmarked on offline accuracy at fixed
query counts.

\paragraph{Our claim.}
This paper asks whether a deterministic, training-free routing layer can
match learned and multi-call routers on accuracy, while preserving an
order-of-magnitude advantage on latency and on per-query cost, and
\emph{scaling cleanly with the agent pool size}. We instantiate this
question as AgentRoute (Section~\ref{sec:design}), an LLM-era adaptation of
the Actor Architecture of~\citet{jang2005efficient}. AgentRoute uses a
substring-based pattern classifier with an LRU cache for the routing
decision, and falls back to a small-model classifier only on low-confidence
queries. The classifier runs in sub-millisecond time, requires no training,
and has zero token cost on cache hits and on confident pattern matches.

\paragraph{Contributions.}
We make four contributions, structured around production-relevant evidence:
\begin{itemize}
    \item A deterministic communication layer for multi-agent LLM systems
    that issues exactly one LLM call per resolved query, with no training and
    no learned components (Section~\ref{sec:design}).
    \item Architecturally faithful re-implementations of five 2025--2026
    routing baselines (MasRouter, RopMura, OI-MAS, RCR-Router, EvoMAS, DyTopo)
    evaluated head-to-head against AgentRoute with a shared
    Qwen 2.5--3B-Instruct backbone~\citep{qwen2025} on
    MedQA-USMLE~\citep{jin2021medqa} and
    GPQA Diamond~\citep{rein2024gpqa} (Section~\ref{sec:benchmarks}).
    \item A scalability study sweeping the registered agent pool from 5 to
    200 agents, characterising how routing-decision latency and per-query
    tokens grow with $N$ for each method (Section~\ref{sec:scalability}).
    \item A production case study against the Claude API with concurrent
    rate-limit handling, budget tracking, and a deployment-narrative
    discussion of what failed and what we learned
    (Section~\ref{sec:casestudy}).
\end{itemize}

We do not claim that pattern-based routing is universally optimal: it
inherits the limitations of its keyword inventory and cannot perform
multi-hop decomposition. Our contribution is to characterise the regime
where the routing layer's fixed cost is itself the binding production
constraint, and to argue that this regime is common enough to warrant a
deterministic baseline alongside the more expressive learned routers.

\section{Related Work}
\label{sec:related}

\paragraph{Multi-agent LLM systems.}
The dominant production paradigm decomposes tasks across role-typed agents
that exchange natural-language messages. AutoGen formalises conversable
multi-agent group chat~\citep{wu2023autogen}; MetaGPT instantiates
SOP-driven assembly lines for software
engineering~\citep{hong2024metagpt}; CAMEL~\citep{li2023camel},
ChatDev~\citep{qian2024chatdev}, and AgentVerse~\citep{chen2024agentverse}
explore role-playing, chat-chain communication, and dynamic team formation
respectively. \citet{cemri2025why} provide an empirical failure taxonomy
across seven frameworks and attribute a substantial fraction of failures
to inter-agent communication errors, motivating selective routing.

\paragraph{LLM routing for single-agent inference.}
A parallel line of work routes individual queries to one of several
available LLM backbones to trade quality for cost.
\citet{chen2023frugalgpt} introduce the prompt-adaptation/cascade
taxonomy. HybridLLM predicts query difficulty to choose between a small
and a large model~\citep{ding2024hybridllm}; RouteLLM learns the router
from preference data~\citep{ong2024routellm}; BEST-Route combines model
selection with test-time best-of-$n$~\citep{ding2025bestroute}. These
methods route \emph{queries to models} rather than \emph{messages between
agents}, but they share the underlying goal of reducing avoidable LLM
calls. AgentRoute applies the same cost-aware principle one layer up, at
the inter-agent communication boundary.

\paragraph{Routing and topology design within multi-agent systems.}
Closer to our setting, GPTSwarm represents an LLM agent system as an
optimisable graph~\citep{zhuge2024gptswarm}; G-Designer uses GNNs to
design task-aware topologies offline~\citep{zhang2025gdesigner}; MaAS
samples query-dependent agent architectures from an agentic
supernet~\citep{zhang2025maas}; AFlow performs MCTS over agentic
workflows~\citep{zhang2025aflow}; AgentPrune prunes communication edges
once and reuses the pruned graph~\citep{zhang2025agentprune}. These
methods optimise the \emph{structure} of multi-agent communication;
AgentRoute is complementary in that it provides a deterministic
\emph{runtime} routing mechanism that can sit on top of any of these
topologies.

\paragraph{Direct baselines.}
The five direct baselines we compare against introduce additional cost in
the routing layer itself. MasRouter trains a cascaded controller and
invokes it per query~\citep{yue2025masrouter}. RopMura's planner uses
four LLM-driven submodules with multi-hop
iteration~\citep{wang2025ropmura}. OI-MAS adds a verifier and an optional
refiner per query~\citep{wang2026oimas}. RCR-Router operates over three
roles (Planner, Searcher, Recommender) for $T$ iterative rounds under
per-role token budgets, with default $T{=}3$~\citep{liu2025rcrrouter}.
EvoMAS evolves a pool of multi-agent configurations and samples one per
query~\citep{evomas2026}. DyTopo runs a Manager plus four agents that
rebuild the directed communication graph each round through embedding
matching~\citep{lu2026dytopo}. AgentRoute occupies the opposite end of
this design spectrum: a deterministic routing decision followed by a
single LLM call to the chosen specialist.

\paragraph{Actor-based middleware.}
AgentRoute is an LLM-era adaptation of the Actor
Architecture~\citep{jang2005efficient}, which introduced location-aware
agent naming, middle-agent matchmaking, and delayed message buffering for
mobile agents. We inherit the first two and replace the third (active
brokering with mobile search objects) with deterministic pattern
matching. This is a deliberate trade-off: we sacrifice the expressiveness
of agent-supplied search code in exchange for predictable routing
latency, an auditable decision trace, and zero per-routing-decision
compute cost---properties that matter more than expressiveness in
production deployments.

% =====================================================================
% Sections 4-7 + Limitations to follow.
% Section 4: Benchmark Evaluation (MedQA)
% Section 5: Scalability Study
% Section 6: Production Case Study (Claude API)
% Section 7: Discussion
% Section 8: Conclusion
% Limitations (mandatory, before references)
% Appendix
% =====================================================================

\section{AgentRoute Design}
\label{sec:design}

AgentRoute is a thin communication layer placed between user-issued queries
and a pool of registered specialist agents. Each query is processed by a
five-step pipeline that issues at most one LLM call for routing and exactly
one LLM call for answering. We describe the pipeline
(\S\ref{sec:design:pipeline}), the design decisions inherited from the Actor
Architecture~\citep{jang2005efficient} (\S\ref{sec:design:actor}), and the
properties of the resulting system (\S\ref{sec:design:properties}).

\subsection{Pipeline}
\label{sec:design:pipeline}

\paragraph{Step 1: Pattern-based domain classification.}
Each domain $d \in \mathcal{D}$ is associated with a hand-curated keyword
inventory $K_d$. For an incoming query $q$ we compute a per-domain score
$s_d(q) = \sum_{k \in K_d} \mathbb{1}[k \in q] \cdot |k|$, in which longer
keywords contribute more, and pick the highest-scoring domain. Confidence is
the top score normalised by the second-best, bounded in $[0, 1]$. The
classifier is implemented as a substring scan and runs in sub-millisecond
time on commodity CPUs; it issues no LLM calls.

\paragraph{Step 2: Confidence gating.}
If confidence $\geq \tau$ (we use $\tau{=}0.6$ throughout), the predicted
domain is accepted. Otherwise, a small-model fallback classifier is invoked
to re-classify $q$ over $\mathcal{D}$. This is the only step in which
AgentRoute spends LLM tokens for routing, and only on queries the pattern
matcher cannot resolve.

\paragraph{Step 3: LRU result cache.}
A bounded-capacity cache mapping query hash to (domain, response) is
consulted before any agent dispatch. On a cache hit, both the routing
decision and the agent's response are reused with zero additional LLM
calls. The cache is most effective for repeated or near-repeated queries;
on workloads composed entirely of distinct queries the hit rate is zero,
which we treat as the adversarial worst case.

\paragraph{Step 4: Specialist selection.}
A domain-indexed registry maps each domain to a set of specialist agents.
Within the chosen domain, the least-loaded agent (by current outstanding
request count) is selected. If every agent in the chosen domain is above
an 85\,\% load threshold, the next-most-relevant domain is tried. Selection
is deterministic given the load state and uses no LLM calls.

\paragraph{Step 5: Single specialist invocation.}
The selected agent receives the original query and produces the response in
a single LLM call. AgentRoute thus issues exactly one LLM call to the
chosen specialist per resolved query, with the optional fallback in Step 2
being the only place additional tokens are consumed for routing.

\subsection{Relation to the Actor Architecture}
\label{sec:design:actor}

AgentRoute instantiates two of the three core ideas
of~\citet{jang2005efficient} in the LLM setting. The first is
location-aware addressing: in the original Actor Architecture, agents carry
both a universal name (UAN) and a location-based name (LAN) that updates as
the agent migrates between hosts; in AgentRoute, the location is the
domain registry entry, and the specialist-selection step plays the role of
the Message Manager's directed forwarding. The second is middle-agent
matchmaking: the Actor Architecture's ATSpace tuple space accepts service
queries from clients and returns matching service agents; AgentRoute's
pattern classifier plus LRU cache plays the same role for natural-language
queries.

The third Actor Architecture idea---\emph{active brokering with mobile
search objects}, in which clients send executable code into the middle
agent's address space to perform expressive searches---is deliberately
not inherited. We replace it with deterministic pattern matching. This
trades the expressiveness of agent-supplied search code for three
properties that matter more in production deployments: predictable
sub-millisecond routing latency, an auditable decision trace (every
routing decision is a deterministic function of the query and the
keyword inventory), and zero per-routing-decision compute cost.

\subsection{Properties}
\label{sec:design:properties}

\paragraph{No training, no learned components.}
The keyword inventory is the only domain-specific component and is edited
by hand. AgentRoute requires no offline training data and has no learned
parameters. This contrasts with MasRouter~\citep{yue2025masrouter},
which trains a cascaded controller, and with methods such as
RopMura~\citep{wang2025ropmura} and RCR-Router~\citep{liu2025rcrrouter}
that rely on dataset-tuned prompts or learned embeddings for their
internal routing decisions.

\paragraph{Constant routing cost in agent-pool size.}
The routing-decision cost is $O(|\mathcal{D}| \cdot \max_d |K_d|)$ for the
pattern scan, plus $O(1)$ amortised for the cache lookup, plus
$O(|\mathcal{A}_d|)$ for the within-domain least-loaded selection where
$\mathcal{A}_d$ is the set of agents registered to the chosen domain. None
of these depend on the total number of registered agents $N$ when the
domain decomposition is held fixed. In contrast, learned controllers and
topology-rebuilding methods score over or embed the candidate agent set,
producing routing cost that grows with $N$.

\paragraph{Auditable decision trace.}
Every routing decision is a deterministic function of $(q, K, \tau, \text{cache state})$
and can be replayed offline without re-invoking an LLM. This matters for
deployments in regulated settings, where the routing layer must be
explainable and reproducible from logs alone.

\paragraph{Trade-off: limited expressiveness.}
The deterministic pattern matcher cannot perform multi-hop decomposition
or recover from a miscalibrated keyword inventory. On benchmarks where
multi-call routers achieve higher accuracy through iterative
decomposition---knowledge-intensive multi-hop question answering, in
particular---AgentRoute will accept a lower accuracy in exchange for the
cost and scalability properties above. We quantify this trade-off
empirically in Sections~\ref{sec:benchmarks} and~\ref{sec:scalability}.

\bibliography{references}

\end{document}
